\begin{notechapter}
\chapter{Empirical Risk Minimization}\label{ch:erm}
\chaptersubtitle{Approximation, estimation and the bias--variance trade-off}

\begin{chapterintro}
The Bayes predictor of \cref{ch:bayes-decision-theory} requires the unknown
distribution $P$. Learning algorithms replace $P$ by the sample and choose,
within a fixed class of models, the one with the smallest training error. This
chapter studies that principle. You will learn
\begin{itemize}[nosep]
  \item the empirical risk minimization (ERM) principle and the role of the
        hypothesis class;
  \item the difference between the risk of \emph{one} trained model and the
        expected risk of a \emph{learning algorithm};
  \item how the excess risk splits into an approximation error and an
        estimation error, and why the two pull in opposite directions;
  \item why the training error is optimistic, with an exact formula for
        linear regression;
  \item the bias--variance decomposition and how regularization trades one
        for the other;
  \item why ERM works when the empirical risk converges to the risk
        \emph{uniformly} over the class, and a first generalization bound for
        finite classes.
\end{itemize}
\prerequisites{\cref{ch:learning-problem,ch:bayes-decision-theory};
Hoeffding's inequality (stated below); basic linear algebra for
\cref{erm:ex:ols-optimism}.}
\end{chapterintro}

\section{The ERM principle}\label{erm:sec:principle}

\begin{definition}[Hypothesis class and ERM]\label{erm:def:erm}
A \emph{hypothesis class} $\Hs$ is a set of predictors $\Xs\to\Ys$ fixed before
seeing the data. The \emph{empirical risk minimizer} over $\Hs$ is
\[
  \hat f_n \in \argmin_{f \in \Hs} \ERisk(f),
  \qquad \ERisk(f) = \frac1n\sum_{i=1}^n \ell\bigl(f(X_i), Y_i\bigr).
\]
\end{definition}

Most familiar training procedures are instances of ERM.

\begin{example}[Least squares and logistic regression]
\emph{Linear regression}: $\Hs = \{x \mapsto w^\top x + b\}$ with the squared loss;
ERM is ordinary least squares, minimizing $\frac1n\sum_i (w^\top x_i + b - y_i)^2$.

\emph{Logistic regression}: $\Hs = \{x \mapsto \sigma(w^\top x + b)\}$ with the
log-loss; ERM is maximum likelihood estimation of $(w,b)$.

\emph{Neural networks}: $\Hs$ is the set of functions computed by a fixed
architecture as its weights vary; training by gradient descent searches (only
approximately) for the empirical risk minimizer.
\end{example}

The hypothesis class encodes what we believe about the solution before seeing
any data --- smoothness, linearity, invariances. It is indispensable: without
it, ERM fails completely.

\begin{pitfall}[The memorizer]\label{erm:pit:memorizer}
Let $X$ be uniform on $[0,1]$, $\Ys=\{0,1\}$, and let $\Hs$ contain \emph{all}
functions. The predictor
\[
  \hat f_n(x) = \begin{cases} y_i & \text{if } x = x_i \text{ for some } i,\\ 0 & \text{otherwise}\end{cases}
\]
has zero training error, so it is an empirical risk minimizer. But a fresh
input almost surely differs from every training input, so
$\Risk(\hat f_n) = \Prob(Y \neq 0)$, which may be as large as $1$. The
empirical risk of $\hat f_n$ says nothing about its risk: $\hat f_n$ has been
tailored to the sample, and \cref{learning-problem:prop:unbiased} does not
apply to it.
\end{pitfall}

\section{The risk of a model and the risk of an algorithm}\label{erm:sec:model-vs-algorithm}

A \emph{learning algorithm} is a map $\mathcal A$ from datasets to predictors,
$\hat f = \mathcal A(D_n)$. Two different quantities measure its quality.

\begin{description}
  \item[Risk of the trained model.] For the observed dataset $d_n$, the model
        $\hat f_{d_n} = \mathcal A(d_n)$ is a fixed function and
        \[
          \Risk(\hat f_{d_n}) = \E_{(X,Y)\sim P}\bigl[\ell(\hat f_{d_n}(X), Y)\bigr]
        \]
        is a number: how well \emph{this} model will do in deployment.
  \item[Expected risk of the algorithm.] Treating the training set as random,
        $\Risk(\hat f_{D_n})$ is a random variable, and its mean
        \[
          \E_{D_n}\bigl[\Risk(\hat f_{D_n})\bigr]
          = \E_{D_n,\,(X,Y)}\bigl[\ell(\hat f_{D_n}(X), Y)\bigr]
        \]
        measures how well the \emph{procedure} $\mathcal A$ does on average
        over training sets of size $n$. This is the quantity to compare when
        choosing between algorithms.
\end{description}

\section{Decomposing the excess risk}\label{erm:sec:decomposition}

Let $f_{\Hs}$ denote the best predictor \emph{in the class},
\[
  f_{\Hs} \in \argmin_{f\in\Hs} \Risk(f)
\]
(if the minimum is not attained, replace $\Risk(f_\Hs)$ by $\inf_{f\in\Hs}\Risk(f)$
below). The gap between the learned model and the Bayes risk splits into two
terms of very different nature:
\begin{equation}\label{erm:eq:decomposition}
  \Risk(\hat f_n) - R^*
  = \underbrace{\Risk(\hat f_n) - \Risk(f_\Hs)}_{\text{estimation error}}
  + \underbrace{\Risk(f_\Hs) - R^*}_{\text{approximation error}} .
\end{equation}

\begin{description}
  \item[Approximation error.] The price of searching in $\Hs$ instead of among
        all functions. It does not depend on the data, only on how well $\Hs$
        can represent the Bayes predictor. \Cref{bayes-decision-theory:ex:parabola}
        computed one: fitting lines to a parabola costs $\frac4{45}$.
  \item[Estimation error.] The price of knowing $P$ only through $n$ examples:
        ERM finds the model that is best \emph{on the sample}, not the one
        that is best \emph{on average}. It is random, non-negative, and
        typically shrinks as $n$ grows.
\end{description}

\begin{figure}[ht]
\centering
\begin{tikzpicture}[>=Stealth, font=\small]
  \draw[thick, fill=gray!6] (0,0) ellipse (5.2cm and 2.3cm);
  \node at (-3.3,2.65) {all predictors};
  \draw[thick, bookaccent, fill=bookaccent!8] (-1.4,-0.2) ellipse (2.4cm and 1.4cm);
  \node[bookaccent] at (-2.7,0.75) {$\Hs$};
  \fill (3.4,0.3) circle (2pt) node[right] {$f^*$ (Bayes)};
  \fill[bookaccent] (0.7,0.1) circle (2pt) node[left] {$f_\Hs$};
  \fill[bookpitfall] (-0.6,-0.7) circle (2pt) node[below] {$\hat f_n$};
  \draw[<->, thick] (0.85,0.12) -- node[above] {approximation} (3.25,0.3);
  \draw[<->, thick, bookpitfall] (-0.5,-0.62) -- node[below right] {estimation} (0.6,0.0);
\end{tikzpicture}
\caption{The excess risk of the learned model $\hat f_n$ decomposes into the
distance from $\Hs$ to the Bayes predictor (approximation) and the distance from
$\hat f_n$ to the best model of the class (estimation). ``Distance'' is meant
in terms of risk.}
\label{erm:fig:decomposition}
\end{figure}

Adding $R^*$ to both sides of \eqref{erm:eq:decomposition} gives
\[
  \Risk(\hat f_n) = \text{estimation error} + \text{approximation error} + \underbrace{R^*}_{\text{Bayes risk}},
\]
and taking expectations over $D_n$ gives the same decomposition for the
expected risk of the algorithm.

\subsection{The fundamental trade-off}

Enlarging $\Hs$ can only decrease the approximation error, since the best model
in a bigger class is at least as good. But it gives ERM more ways to fit the
noise of the sample, so it typically increases the estimation error.
\Cref{erm:fig:tradeoff} shows the resulting picture: the total risk is
U-shaped in the complexity of the class.
\begin{itemize}
  \item A class that is too small \emph{underfits}: large approximation error,
        and the training error itself stays high.
  \item A class that is too large \emph{overfits}: small training error, large
        estimation error, large gap between training and test error.
\end{itemize}

\begin{figure}[ht]
\centering
\begin{tikzpicture}
\begin{axis}[width=0.55\textwidth, height=5.4cm, xmin=0.3, xmax=6, ymin=0, ymax=2.2,
  xlabel={complexity of $\Hs$}, ylabel={risk}, xtick=\empty, ytick=\empty,
  axis lines=left, legend style={at={(1.03,0.5)}, anchor=west, font=\small},
  samples=80]
  \addplot[thick, bookdef, domain=0.3:6] {1.2/x};
  \addlegendentry{approximation}
  \addplot[thick, bookpitfall, domain=0.3:6] {0.05*x^1.6};
  \addlegendentry{estimation}
  \addplot[very thick, black, domain=0.3:6] {1.2/x + 0.05*x^1.6 + 0.15};
  \addlegendentry{total risk}
  \addplot[dashed, gray, domain=0.3:6] {0.15};
  \addlegendentry{Bayes risk $R^*$}
  \draw[dotted] (axis cs:2.75,0) -- (axis cs:2.75,1.06);
  \node[font=\small, anchor=north] at (axis cs:1.3,2.1) {underfitting};
  \node[font=\small, anchor=north] at (axis cs:5.0,2.1) {overfitting};
\end{axis}
\end{tikzpicture}
\caption{Schematic trade-off between approximation and estimation error. The
best class complexity (dotted line) balances the two.}
\label{erm:fig:tradeoff}
\end{figure}

\section{The optimism of the training error}\label{erm:sec:optimism}

The training error of the ERM solution is a biased estimate of its risk, and
the bias always goes the same way.

\begin{proposition}[Training error is optimistic]\label{erm:prop:optimism}
If $f_\Hs$ exists,
\[
  \E\bigl[\ERisk(\hat f_n)\bigr] \;\leq\; \Risk(f_\Hs) \;\leq\; \E\bigl[\Risk(\hat f_n)\bigr].
\]
\end{proposition}

\begin{proof}
Since $\hat f_n$ minimizes the empirical risk over $\Hs$ and $f_\Hs\in\Hs$,
$\ERisk(\hat f_n) \leq \ERisk(f_\Hs)$ for every sample. The predictor
$f_\Hs$ does not depend on the sample, so by \cref{learning-problem:prop:unbiased}
$\E[\ERisk(f_\Hs)] = \Risk(f_\Hs)$, which gives the first inequality. The second
holds because $\Risk(\hat f_n) \geq \Risk(f_\Hs)$ for every sample, by definition
of $f_\Hs$.
\end{proof}

The expected training error is below the best achievable risk, and the
expected risk is above it: the training error underestimates the risk of the
model it produced. For linear regression the gap can be computed exactly.

\begin{example}[Optimism of least squares]\label{erm:ex:ols-optimism}
Consider a fixed design matrix $\mathbf X \in \R^{n\times p}$ of full rank
$p$ and responses $y = \mathbf X\beta + \varepsilon$, where the noise vector
$\varepsilon$ has independent entries with mean $0$ and variance $\sigma^2$. The
least-squares fit is $\hat y = H y$, with $H = \mathbf X(\mathbf X^\top\mathbf X)^{-1}\mathbf X^\top$
the orthogonal projection onto the column space of $\mathbf X$; $H$ is
symmetric, idempotent, and $\tr H = \tr\bigl((\mathbf X^\top\mathbf X)^{-1}\mathbf X^\top\mathbf X\bigr) = p$.

\emph{Training error.} Since $(I - H)\mathbf X\beta = 0$, the residual is
$y - \hat y = (I-H)\varepsilon$, and for any matrix $A$,
$\E\norm{A\varepsilon}^2 = \sigma^2 \tr(A^\top A)$. Hence
\[
  \E\Bigl[\tfrac1n\norm{y - \hat y}^2\Bigr] = \tfrac{\sigma^2}{n}\tr(I-H) = \sigma^2\Bigl(1 - \frac pn\Bigr).
\]
\emph{Error on new responses.} Let $y' = \mathbf X\beta + \varepsilon'$ be new
responses at the same inputs, with noise $\varepsilon'$ independent of
$\varepsilon$. Then $y' - \hat y = \varepsilon' - H\varepsilon$ and
\[
  \E\Bigl[\tfrac1n\norm{y' - \hat y}^2\Bigr] = \tfrac{1}{n}\bigl(n\sigma^2 + \sigma^2 \tr H\bigr) = \sigma^2\Bigl(1 + \frac pn\Bigr).
\]
The Bayes risk is $\sigma^2$. The training error is below it by
$\sigma^2 p/n$ and the true error above it by the same amount: the optimism is
$2\sigma^2 p/n$, proportional to the number of parameters and inversely
proportional to the sample size. With $p = n$ the training error is zero ---
the model interpolates the data.
\end{example}

This is why model performance is always reported on data not used for
training. A \emph{test set} $\{(x'_j,y'_j)\}_{j=1}^m$ drawn independently of
the training set gives
\[
  \widehat R_{\mathrm{test}}(\hat f) = \frac1m\sum_{j=1}^m \ell\bigl(\hat f(x'_j), y'_j\bigr),
\]
which, conditionally on the trained model, is an unbiased estimate of
$\Risk(\hat f)$ with variance $\Var(\ell(\hat f(X),Y))/m$.

\begin{pitfall}[A test set used for selection is no longer a test set]
If the test set is used to choose among models (or hyperparameters), the
selected model has been fitted to the test set too, and its test error is
optimistic for the same reason as the training error. The usual remedy is a
three-way split: a training set to fit, a \emph{validation} set to select, and
a test set used once at the very end. \Cref{ch:measuring-performance} describes
cross-validation, which reuses the data more efficiently.
\end{pitfall}

\section{Bias, variance and regularization}\label{erm:sec:bias-variance}

For the squared loss the estimation error has a classical finer structure.
Fix an input $x$ and consider the prediction $\hat f_{D_n}(x)$ as a random
variable over training sets, with mean $\bar f(x) = \E_{D_n}[\hat f_{D_n}(x)]$.

\begin{proposition}[Bias--variance decomposition]\label{erm:prop:bias-variance}
Let $m(x) = \E[Y \mid X=x]$. For a test input $X$ independent of the training
set,
\[
  \E_{D_n}\bigl[\Risk(\hat f_{D_n})\bigr]
  = \underbrace{\E\bigl[\Var(Y\mid X)\bigr]}_{\text{noise } = R^*}
  + \underbrace{\E\bigl[(\bar f(X) - m(X))^2\bigr]}_{\text{squared bias}}
  + \underbrace{\E\bigl[\Var_{D_n}(\hat f_{D_n}(X))\bigr]}_{\text{variance}} .
\]
\end{proposition}

\begin{proof}
By \eqref{bayes-decision-theory:eq:pythagoras} applied to the fixed function
$\hat f_{D_n}$, $\Risk(\hat f_{D_n}) = R^* + \E_X[(\hat f_{D_n}(X) - m(X))^2]$.
Take the expectation over $D_n$ and, for each fixed $x$, decompose
$\E_{D_n}[(\hat f_{D_n}(x) - m(x))^2]$ as for any random variable $Z$ and constant
$c$: $\E[(Z - c)^2] = (\E Z - c)^2 + \Var Z$, with $Z = \hat f_{D_n}(x)$ and $c = m(x)$.
\end{proof}

The \emph{bias} measures how far the average model is from the truth --- it is
large for a rigid class (the horizontal line of
\cref{bayes-decision-theory:ex:parabola} has a large bias). The
\emph{variance} measures how much the model changes from one training set to
another --- it is large for flexible classes and small samples.

\subsection{Regularization}

Instead of choosing a smaller class, one can keep a rich class and add a
penalty $\Omega(f)$ that discourages complex solutions:
\[
  \hat f_{n,\lambda} \in \argmin_{f\in\Hs}\ \Bigl\{\ERisk(f) + \lambda\,\Omega(f)\Bigr\},
  \qquad \lambda \geq 0 .
\]
The strength $\lambda$ moves the solution along the trade-off of
\cref{erm:fig:tradeoff}: $\lambda = 0$ is plain ERM, and $\lambda \to \infty$
forces the simplest model.

\begin{example}[Ridge regression]
For linear models $f(x) = w^\top x$ with the squared loss and
$\Omega(w) = \norm{w}^2$, the objective
$\frac1n\norm{y - \mathbf Xw}^2 + \lambda\norm{w}^2$ is a strictly convex
quadratic when $\lambda > 0$. Setting its gradient
$-\frac2n\mathbf X^\top(y - \mathbf Xw) + 2\lambda w$ to zero gives
\[
  \hat w_\lambda = \bigl(\mathbf X^\top\mathbf X + n\lambda I\bigr)^{-1}\mathbf X^\top y .
\]
The added $n\lambda I$ makes the matrix invertible even with collinear
features, and shrinks the coefficients toward zero: the estimate is biased,
but its variance is smaller. With $\Omega(w) = \norm{w}_1$ (the \emph{lasso})
some coefficients become exactly zero, which also performs variable selection.
\end{example}

In practice $\lambda$ (like the choice of $\Hs$) is a \emph{hyperparameter},
chosen by evaluating several values on validation data.

\section{When does ERM work? Uniform convergence}\label{erm:sec:uniform}

We now look for guarantees that the estimation error is small.

\begin{definition}[Consistency]
ERM over $\Hs$ is \emph{consistent} for $P$ if its estimation error vanishes:
$\Risk(\hat f_n) \to \Risk(f_\Hs)$ in probability as $n\to\infty$.
\end{definition}

The law of large numbers does not settle the question. It gives
$\ERisk(f) \to \Risk(f)$ for each \emph{fixed} $f$, but $\hat f_n$ is not fixed:
it is chosen after seeing the sample, precisely because its empirical risk is
small (recall the memorizer of \cref{erm:sec:principle}). What we need is that the empirical risk be
close to the risk \emph{simultaneously for all} $f \in \Hs$; then it cannot
mislead ERM about any model, including the one it selects.

\begin{theorem}[Estimation error and uniform deviation]\label{erm:thm:uniform}
For every sample,
\[
  \Risk(\hat f_n) - \Risk(f_\Hs) \;\leq\; 2\,\sup_{f\in\Hs}\,\abs{\ERisk(f) - \Risk(f)} .
\]
\end{theorem}

\begin{proof}
Insert the empirical risks and use $\ERisk(\hat f_n) \leq \ERisk(f_\Hs)$:
\[
  \Risk(\hat f_n) - \Risk(f_\Hs)
  = \underbrace{\Risk(\hat f_n) - \ERisk(\hat f_n)}_{\leq\ \sup_f \abs{\ERisk(f) - \Risk(f)}}
  + \underbrace{\ERisk(\hat f_n) - \ERisk(f_\Hs)}_{\leq\ 0}
  + \underbrace{\ERisk(f_\Hs) - \Risk(f_\Hs)}_{\leq\ \sup_f\abs{\ERisk(f) - \Risk(f)}} . \qedhere
\]
\end{proof}

\begin{keyidea}[Uniform convergence]
If $\sup_{f\in\Hs}\abs{\ERisk(f) - \Risk(f)} \to 0$ in probability, ERM is
consistent. Proving such \emph{uniform convergence}, with explicit rates, is
the main business of statistical learning theory. Whether it holds depends on
the ``size'' of $\Hs$: it fails for the class of all functions (the
memorizer) and holds for classes of finite VC dimension
(\cref{ch:symmetrization-lemma}).
\end{keyidea}

\subsection{A first bound: finite classes}

For a single fixed predictor, a concentration inequality says how fast the
empirical risk approaches the risk. Abu-Mostafa's \emph{Learning From Data}
calls this step ``probability to the rescue''.

\begin{lemma}[Hoeffding's inequality]
Let $Z_1,\dots,Z_n$ be independent random variables with values in $[a,b]$ and
$\bar Z = \frac1n\sum_i Z_i$. For every $\epsilon > 0$,
\[
  \Prob\bigl(\bar Z - \E\bar Z \geq \epsilon\bigr) \leq e^{-2n\epsilon^2/(b-a)^2},
  \qquad
  \Prob\bigl(\abs{\bar Z - \E\bar Z} \geq \epsilon\bigr) \leq 2\,e^{-2n\epsilon^2/(b-a)^2}.
\]
\end{lemma}

With many hypotheses, the probability that \emph{at least one} of them is
misleading is at most the sum of the individual probabilities (the
\emph{union bound}).

\begin{theorem}[Uniform convergence for a finite class]\label{erm:thm:finite}
Let $\Hs$ be finite and the loss take values in $[0,1]$. For every
$\epsilon>0$,
\[
  \Prob\Bigl(\sup_{f\in\Hs}\abs{\ERisk(f) - \Risk(f)} \geq \epsilon\Bigr) \leq 2\,\abs{\Hs}\,e^{-2n\epsilon^2}.
\]
Equivalently, for $\delta\in(0,1)$, with probability at least $1-\delta$,
\[
  \Risk(\hat f_n) - \Risk(f_\Hs) \;\leq\; 2\sqrt{\frac{\log(2\abs{\Hs}/\delta)}{2n}} .
\]
\end{theorem}

\begin{proof}
The event that $\sup_f\abs{\ERisk(f) - \Risk(f)}\geq\epsilon$ is the union,
over $f\in\Hs$, of the events $\abs{\ERisk(f) - \Risk(f)}\geq\epsilon$. Each $f$ is
fixed, so Hoeffding's inequality applies to the losses
$\ell(f(X_i),Y_i)\in[0,1]$ and bounds the probability of each event by
$2e^{-2n\epsilon^2}$; the union bound adds $\abs{\Hs}$ of them. For the second
form, set $2\abs{\Hs}e^{-2n\epsilon^2} = \delta$, solve for $\epsilon$, and apply
\cref{erm:thm:uniform}.
\end{proof}

\begin{example}[Reading the bound]
With $\abs{\Hs} = 1000$ models, $n = 10\,000$ examples and $\delta = 0.05$,
$\sqrt{\log(40\,000)/20\,000}\approx 0.023$: with probability $95\%$ the ERM
solution is within $0.046$ of the best model of the class. The size of the
class enters only through $\log\abs{\Hs}$, so a million models instead of a
thousand would merely double the logarithm.
\end{example}

Most classes of interest --- linear classifiers, neural networks --- are
infinite, and the union bound over $\abs{\Hs}$ terms is useless. The remedy,
developed in \cref{ch:symmetrization-lemma}, is to notice that on a finite
sample an infinite class behaves like a finite one: only the \emph{patterns}
of predictions it can produce on the sample matter.

\begin{chaptersummary}
\begin{itemize}[nosep]
  \item ERM minimizes the training error over a hypothesis class $\Hs$ fixed in
        advance; without such a restriction it can memorize the sample.
  \item The excess risk $\Risk(\hat f_n) - R^*$ is the sum of an estimation
        error (finite sample) and an approximation error (restricted class);
        richer classes trade the second for the first.
  \item The training error of the ERM solution is optimistic; for least
        squares the optimism is exactly $2\sigma^2p/n$. Risk is estimated on
        independent test data.
  \item For the squared loss, the expected risk is noise $+$ bias$^2$ $+$
        variance; regularization reduces variance at the cost of bias.
  \item The estimation error is at most twice the uniform deviation
        $\sup_f\abs{\ERisk(f) - \Risk(f)}$; for a finite class Hoeffding's
        inequality and the union bound give a rate
        $\sqrt{\log\abs{\Hs}/n}$.
\end{itemize}
\end{chaptersummary}

\section*{Exercises}
\addcontentsline{toc}{section}{Exercises}

\begin{exercise}
Show that the approximation error is a non-increasing function of the class:
if $\Hs_1\subseteq\Hs_2$ then $\Risk(f_{\Hs_2}) \leq \Risk(f_{\Hs_1})$. Is the
estimation error necessarily non-decreasing? Give an example or an argument.
\end{exercise}

\begin{exercise}
In \cref{erm:ex:ols-optimism}, show that $\frac{1}{n-p}\norm{y-\hat y}^2$ is an
unbiased estimator of $\sigma^2$. Why does one divide by $n-p$ rather than $n$?
\end{exercise}

\begin{exercise}
Let $\hat f(x) = \frac1n\sum_i Y_i$ be the constant predictor equal to the
sample mean, with $\Var(Y) = v$. Compute its bias and variance at any $x$ and
its expected risk under the squared loss.
\hint{$\Risk(c) = \Var(Y) + (c - \E Y)^2$ for a constant $c$.}
\end{exercise}

\begin{exercise}
How large must $n$ be in \cref{erm:thm:finite} to guarantee, with
probability $0.99$, an estimation error below $0.01$ for a class of $10^6$
models?
\end{exercise}
\end{notechapter}
